% ---
% Subsección nueva en Docs/Technical_Report.tex
% Decisiones: D-033, H-038, H-039
% ---

\subsubsection{Dimensión efectiva}
\label{sec:dimension_efectiva}

Para caracterizar la expresividad y la estructura espectral de la matriz de kernel sin
depender de las etiquetas de clase, Huang et al. \cite{huang2021power} propusieron el concepto
de \textit{dimensión efectiva} $d(K)$. Esta métrica cuantifica cuántas direcciones
independientes u ortogonales explora el conjunto de datos en el espacio de Hilbert inducido
por el kernel.

Sea $K \in \mathbb{R}^{N \times N}$ una matriz de kernel normalizada a traza $N$
($\operatorname{Tr}(K) = N$). Sean $t_1 \ge t_2 \ge \dots \ge t_N \ge 0$ los autovalores
ordenados de forma decreciente de $K$. La dimensión efectiva se define como:
\begin{equation}
  d(K) = \sum_{k=1}^{N} \frac{1}{N - k + 1}\sum_{l=k}^{N} t_l.
  \label{eq:dim_efectiva}
\end{equation}
En la ecuación~\eqref{eq:dim_efectiva}, el término interior $\sum_{l=k}^{N} t_l$ representa
la suma de la cola espectral a partir del $k$-ésimo autovalor, y el denominador
la promedia entre las $N-k+1$ direcciones restantes.

% respaldo: «In the numerical experiment, we choose the following measure as the approximate dimension» (huang2021powerSupp, p. 12)
% La errata está igual en la versión publicada (S64) y en la de arXiv (F12), sección 6 B = Apéndice F.2.
\textbf{Nota sobre errata en la literatura:} La ecuación (S64) de la información suplementaria de Huang et al. \cite{huang2021power} (sección~6\,B) imprime el factor normalizador como $\frac{1}{N - k}$, formulación que incurre en una indeterminación por división entre cero en el último sumando ($k = N$). La corrección consiste en emplear el denominador $N - k + 1$. Con esta formulación, la métrica satisface las cotas que enuncian los autores:
\begin{equation}
  1 \le d(K) \le N.
  \label{eq:effective_dim_bounds}
\end{equation}
Los dos extremos se comprueban directamente:
\begin{itemize}
  \item \textbf{Cota inferior ($d = 1$):} Si todas las muestras mapean al mismo estado cuántico, la matriz de kernel tiene rango 1 ($t_1 = N$ y $t_l = 0$ para $l \ge 2$). La suma total colapsa a $d = \frac{1}{N} t_1 = \frac{N}{N} = 1$.
  \item \textbf{Cota superior ($d = N$):} Si todas las muestras mapean a estados mutuamente ortogonales en el espacio de Hilbert ($K = I$, autovalores $t_l = 1$ para todo $l$), la cola espectral para cada índice $k$ satisface $\sum_{l=k}^N 1 = N - k + 1$. Al cancelarse con el denominador, cada sumando vale idénticamente 1, resultando en $d = \sum_{k=1}^N 1 = N$.
\end{itemize}

Valores de $d$ cercanos a $N$ indican que cada lesión ocupa casi su propia dirección, como ocurre con un kernel concentrado; valores pequeños, que el kernel agrupa las lesiones en pocas direcciones.
